\thispagestyle{empty}

\begin{center}
    \textbf{\large ABSTRACT}
\end{center}

\vspace{14pt}

\begin{singlespace}
\noindent\justifying
This paper presents the development of WeFIND, a collaborative mobile application designed to connect pet owners and the community in the search, rescue, and dissemination of information about lost and found domestic animals, as well as to encourage responsible adoption. The project was based on field research with pet owners and a comparative analysis of similar platforms, identifying difficulties in existing search methods, such as dependence on street posters and social media groups, and structuring software requirements through UML modeling diagrams. The proposal integrates geolocation features for displaying reports on an interactive map, searches based on the animal's physical characteristics, an in-app private messaging channel that avoids exposing personal phone numbers, and standardized poster generation with QR codes for sharing or printing. Technological development uses React Native and Expo in the mobile environment and the Supabase platform with a PostgreSQL relational database for authentication, storage, and data synchronization services. As a contribution, this work delivers a functional solution to support community mobilization in animal rescue. Exploratory manual functional checks and a usability evaluation with end users indicated that the flows covered in the evaluation operated as expected and that participants perceived the application as easy to use, considering the limitations of the sample and procedures.

\vspace{18pt}

\noindent
\textbf{Keywords:} Lost animals. Mobile applications. Geolocation. React Native. Supabase. Collaborative systems.
\end{singlespace}

\clearpage
